\documentclass{article}
	\usepackage[UTF8]{ctex}
	\usepackage{amsmath}
	\usepackage{graphicx}
	\usepackage{underscore}
	\usepackage{geometry}
	\usepackage{anysize}
	\date{}
	\title{成绩分析题解}
	\geometry{top=10mm,bottom=15mm,left=25mm,right=25mm}
	
	\begin{document}
		\maketitle \vspace{-50pt}
		
		\section{模拟}
			\indent 涉及循环，判断，多方面考察。\\
			\indent 使用数组和不使用数组都可以，并没有卡空间。个人推荐使用数组，因为容易写。当遇到数据数目很大的题目时，需考虑不使用数组(这道题就可以不用数组)，或者使用其它存储方式，如桶排存储[0,100]的数。\\
			\indent 注意中位数有可能有两个。不能对两个中位数分别除以2，有可能它们均为奇数。\\
			\indent 题目给出''结果保证为非负整数''和''成绩按照从小到大的顺序排列''，用于降低难度。
			
	\end
	{document}